% GENERATED FILE. Edit canonical lessons, then run npm run generate:books.
% canonical-source-hash: fnv1a64:07f4e4fb4c2adff4
% canonical-lessons: BN-C507-chhidro, BN-C507-plastar, BN-C507-atha, BN-C507-lepa, BN-C507-thoka, BN-R507-review-500-507

\chapter{Cracks and Repairs}
\label{ch:bn-cracks-and-repairs}
\hlchaptermodality{507}

\begin{chapteropening}
\textbf{By the end of this chapter:} \emph{I can talk about cracks, holes, glue and simple repairs.}

Complete the last lesson of chapter 507: I can talk about cracks, holes, glue and simple repairs.
\end{chapteropening}

\section[\texorpdfstring{\bn{ছিদ্র} (chhidrô)}{ছিদ্র (chhidrô)}]{\bn{ছিদ্র} (chhidrô) — a hole; a small opening}

\label{lesson:BN-C507-chhidro}



\begin{quote}
Before the new one: say the Bengali for to bail out water, then the Bengali for to seep.
\end{quote}

\subsection*{You'll want to know: \bn{ছিদ্র}}
\textbf{\bn{ছিদ্র}} — \emph{chhidrô} — \textquotedblleft{}a hole; a small opening\textquotedblright{}.

\begin{quote}
\bn{নৌকার} \bn{তলায়} \bn{একটা} \bn{ছিদ্র} \bn{ছিল।}
\end{quote}

\begin{quote}
\emph{Noukār tôlāy ekṭā chhidrô chhilô.}
\end{quote}

\begin{quote}
There was a hole in the bottom of the boat.
\end{quote}

\begin{cousinweb}
\textbf{\bn{ছিদ্র}} (\emph{chhidrô}) is a tatsama, Sanskrit \emph{chidra}, \textquotedblleft{}torn asunder\textquotedblright{}, and so \textquotedblleft{}a hole, an opening\textquotedblright{} (Monier-Williams).
\end{cousinweb}

\begin{grammarlens}[title={What you've built}]
One more word for describing what you lived and felt.
\end{grammarlens}

\subsection*{Your turn}
\begin{itemize}
\raggedright
  \item \emph{Say it:} \emph{chhidrô}
  \item \emph{Say it:} \emph{chhidrô}, once more
  \item \emph{Say it:} \emph{Noukār tôlāy ekṭā chhidrô chhilô.} — There was a hole in the bottom of the boat.
  \item \emph{Recall:} say the Bengali for to bail out water, then the Bengali for to seep, then say \emph{chhidrô} again
\end{itemize}

\begin{tcolorbox}[breakable,colback=teal!4,colframe=teal!35!black,title={Before you move on}]
What does \bn{ছিদ্র} mean? (\textquotedblleft{}A hole; a small opening\textquotedblright{}.) Say it once more.
\end{tcolorbox}
\section[\texorpdfstring{\bn{প্লাস্টার} (plāsṭār)}{প্লাস্টার (plāsṭār)}]{\bn{প্লাস্টার} (plāsṭār) — wall plaster}

\label{lesson:BN-C507-plastar}



\begin{quote}
Before the new one: say the Bengali for to seep, then the Bengali for a small opening.
\end{quote}

\subsection*{You'll want to know: \bn{প্লাস্টার}}
\textbf{\bn{প্লাস্টার}} — \emph{plāsṭār} — \textquotedblleft{}wall plaster\textquotedblright{}.

\begin{quote}
\bn{পুরনো} \bn{দেওয়াল} \bn{থেকে} \bn{প্লাস্টার} \bn{ঝরছিল।}
\end{quote}

\begin{quote}
\emph{Purôno deoyāl theke plāsṭār jhôrchhilô.}
\end{quote}

\begin{quote}
Plaster was flaking off the old wall.
\end{quote}

\begin{cousinweb}
\textbf{\bn{প্লাস্টার}} (\emph{plāsṭār}) is borrowed from English \emph{plaster}.
\end{cousinweb}

\begin{grammarlens}[title={What you've built}]
One more word for describing what you lived and felt.
\end{grammarlens}

\subsection*{Your turn}
\begin{itemize}
\raggedright
  \item \emph{Say it:} \emph{plāsṭār}
  \item \emph{Say it:} \emph{plāsṭār}, once more
  \item \emph{Say it:} \emph{Purôno deoyāl theke plāsṭār jhôrchhilô.} — Plaster was flaking off the old wall.
  \item \emph{Recall:} say the Bengali for to seep, then the Bengali for a small opening, then say \emph{plāsṭār} again
\end{itemize}

\begin{tcolorbox}[breakable,colback=teal!4,colframe=teal!35!black,title={Before you move on}]
What does \bn{প্লাস্টার} mean? (\textquotedblleft{}Wall plaster\textquotedblright{}.) Say it once more.
\end{tcolorbox}
\section[\texorpdfstring{\bn{আঠা} (āṭhā)}{আঠা (āṭhā)}]{\bn{আঠা} (āṭhā) — glue; paste}

\label{lesson:BN-C507-atha}



\begin{quote}
Before the new one: say the Bengali for a small opening, then the Bengali for wall plaster.
\end{quote}

\subsection*{You'll want to know: \bn{আঠা}}
\textbf{\bn{আঠা}} — \emph{āṭhā} — \textquotedblleft{}glue; paste\textquotedblright{}.

\begin{quote}
\bn{ছেঁড়া} \bn{বইটা} \bn{আঠা} \bn{লাগিয়ে} \bn{জুড়লাম।}
\end{quote}

\begin{quote}
\emph{Chhẽṛā bôiṭā āṭhā lāgiye juṛlām.}
\end{quote}

\begin{quote}
I stuck the torn book back together with glue.
\end{quote}

\begin{cousinweb}
\textbf{\bn{আঠা}} (\emph{āṭhā}) comes from a reconstructed \emph{aṭṭha}; CDIAL lists Bengali \emph{āṭhā}, \textquotedblleft{}gum, paste\textquotedblright{} (CDIAL 186).
\end{cousinweb}

\begin{grammarlens}[title={What you've built}]
One more word for describing what you lived and felt.
\end{grammarlens}

\subsection*{Your turn}
\begin{itemize}
\raggedright
  \item \emph{Say it:} \emph{āṭhā}
  \item \emph{Say it:} \emph{āṭhā}, once more
  \item \emph{Say it:} \emph{Chhẽṛā bôiṭā āṭhā lāgiye juṛlām.} — I stuck the torn book back together with glue.
  \item \emph{Recall:} say the Bengali for a small opening, then the Bengali for wall plaster, then say \emph{āṭhā} again
\end{itemize}

\begin{tcolorbox}[breakable,colback=teal!4,colframe=teal!35!black,title={Before you move on}]
What does \bn{আঠা} mean? (\textquotedblleft{}Glue; paste\textquotedblright{}.) Say it once more.
\end{tcolorbox}
\section[\texorpdfstring{\bn{লেপা} (lepā)}{লেপা (lepā)}]{\bn{লেপা} (lepā) — to plaster; to smear}

\label{lesson:BN-C507-lepa}



\begin{quote}
Before the new one: say the Bengali for wall plaster, then the Bengali for glue.
\end{quote}

\subsection*{You'll want to know: \bn{লেপা}}
\textbf{\bn{লেপা}} — \emph{lepā} — \textquotedblleft{}to plaster; to smear\textquotedblright{}.

\begin{quote}
\bn{মজুর} \bn{কাদা} \bn{লেপে} \bn{দেওয়ালের} \bn{ফুটোটা} \bn{বন্ধ} \bn{করল।}
\end{quote}

\begin{quote}
\emph{Mojur kādā lepe deoyāler phuṭoṭā bôndhô kôrlô.}
\end{quote}

\begin{quote}
The labourer plastered mud over the hole in the wall and closed it.
\end{quote}

Past (he/she …): \bn{তিনি} \bn{দেওয়ালে} \bn{কাদা} \bn{লেপলেন}, \emph{tini deoyāle kādā leplen}.

\begin{cousinweb}
\textbf{\bn{লেপা}} (\emph{lepā}) is a tadbhava: CDIAL lists Bengali \emph{lepā}, \textquotedblleft{}to plaster\textquotedblright{}, under Sanskrit \emph{lepya}, \textquotedblleft{}to be smeared\textquotedblright{} (CDIAL 11114).
\end{cousinweb}

\begin{grammarlens}[title={What you've built}]
One more word for describing what you lived and felt.
\end{grammarlens}

\subsection*{Your turn}
\begin{itemize}
\raggedright
  \item \emph{Say it:} \emph{lepā}
  \item \emph{Say it:} \emph{lepā}, once more
  \item \emph{Say it:} \emph{Mojur kādā lepe deoyāler phuṭoṭā bôndhô kôrlô.} — The labourer plastered mud over the hole in the wall and closed it.
  \item \emph{Recall:} say the Bengali for wall plaster, then the Bengali for glue, then say \emph{lepā} again
\end{itemize}

\begin{tcolorbox}[breakable,colback=teal!4,colframe=teal!35!black,title={Before you move on}]
What does \bn{লেপা} mean? (\textquotedblleft{}To plaster; to smear\textquotedblright{}.) Say it once more.
\end{tcolorbox}
\section[\texorpdfstring{\bn{ঠোকা} (ṭhokā)}{ঠোকা (ṭhokā)}]{\bn{ঠোকা} (ṭhokā) — to hammer in; to knock}

\label{lesson:BN-C507-thoka}



\begin{quote}
Before the new one: say the Bengali for glue, then the Bengali for to plaster.
\end{quote}

\subsection*{You'll want to know: \bn{ঠোকা}}
\textbf{\bn{ঠোকা}} — \emph{ṭhokā} — \textquotedblleft{}to hammer in; to knock\textquotedblright{}.

\begin{quote}
\bn{মিস্ত্রি} \bn{দেওয়ালে} \bn{একটা} \bn{পেরেক} \bn{ঠুকল।}
\end{quote}

\begin{quote}
\emph{Mistri deoyāle ekṭā perek ṭhuklô.}
\end{quote}

\begin{quote}
The mechanic hammered a nail into the wall.
\end{quote}

Past (he/she …): \bn{তিনি} \bn{দরজায়} \bn{ঠুকলেন}, \emph{tini dôrjāy ṭhuklen}.

The stem vowel rises in the past: \bn{ঠুকল}, he hammered.

\begin{cousinweb}
\textbf{\bn{ঠোকা}} (\emph{ṭhokā}) comes from a reconstructed root \emph{ṭhōkk}, \textquotedblleft{}knock\textquotedblright{}; CDIAL lists Bengali \emph{ṭhokā}, \textquotedblleft{}to strike\textquotedblright{} (CDIAL 5513).
\end{cousinweb}

\begin{grammarlens}[title={What you've built}]
One more word for describing what you lived and felt.
\end{grammarlens}

\subsection*{Your turn}
\begin{itemize}
\raggedright
  \item \emph{Say it:} \emph{ṭhokā}
  \item \emph{Say it:} \emph{ṭhokā}, once more
  \item \emph{Say it:} \emph{Mistri deoyāle ekṭā perek ṭhuklô.} — The mechanic hammered a nail into the wall.
  \item \emph{Recall:} say the Bengali for glue, then the Bengali for to plaster, then say \emph{ṭhokā} again
\end{itemize}

\begin{tcolorbox}[breakable,colback=teal!4,colframe=teal!35!black,title={Before you move on}]
What does \bn{ঠোকা} mean? (\textquotedblleft{}To hammer in; to knock\textquotedblright{}.) Say it once more.
\end{tcolorbox}
\section[(dialogue)]{Looking back: from Orders and Departments to Cracks and Repairs}

\label{lesson:BN-R507-review-500-507}



\begin{quote}
Say the Bengali for a department and a command.
\end{quote}

\subsection*{Your turn}
Say these aloud:

\begin{itemize}
\raggedright
  \item a department, a command, a workplace, to get done, to supervise
  \item a specialist, an adviser, a translator, an architect, to diagnose
  \item a computer, a mobile phone, a selfie, to rummage, to charge a phone
  \item a website, a link online, online, a screen, to download
  \item a robot, automatic, a generator, to rust, a motor
  \item darkness, bright, a chandelier, to flicker, to go out (of a light)
  \item a water tank, a hole right through, to overflow, to bail out water, to seep
  \item a small opening, wall plaster, glue, to plaster, to hammer in
\end{itemize}

\begin{tcolorbox}[breakable,colback=teal!4,colframe=teal!35!black,title={Before you move on}]
A department? (\textbf{\bn{বিভাগ}}, \emph{bibhāg}.) A command? (\textbf{\bn{হুকুম}}, \emph{hukum}.) A workplace? (\textbf{\bn{কর্মস্থল}}, \emph{kôrmôsthôl}.) To get done? (\textbf{\bn{সারা}}, \emph{sārā}.) To supervise? (\textbf{\bn{তদারক} \bn{করা}}, \emph{tôdārôk kôrā}.) A specialist? (\textbf{\bn{বিশেষজ্ঞ}}, \emph{bisheshôgyô}.) An adviser? (\textbf{\bn{উপদেষ্টা}}, \emph{upôdeshṭā}.) A translator? (\textbf{\bn{অনুবাদক}}, \emph{ônubādôk}.) An architect? (\textbf{\bn{স্থপতি}}, \emph{sthôpôti}.) To diagnose? (\textbf{\bn{নির্ণয়} \bn{করা}}, \emph{nirnôy kôrā}.) A computer? (\textbf{\bn{কম্পিউটার}}, \emph{kômpiuṭār}.) A mobile phone? (\textbf{\bn{মোবাইল}}, \emph{mobāil}.) A selfie? (\textbf{\bn{সেলফি}}, \emph{selphi}.) To rummage? (\textbf{\bn{ঘাঁটা}}, \emph{ghā̃ṭā}.) To charge a phone? (\textbf{\bn{চার্জ} \bn{দেওয়া}}, \emph{chārj deoyā}.) A website? (\textbf{\bn{ওয়েবসাইট}}, \emph{oyebsāiṭ}.) A link online? (\textbf{\bn{লিংক}}, \emph{lingk}.) Online? (\textbf{\bn{অনলাইন}}, \emph{ônlāin}.) A screen? (\textbf{\bn{স্ক্রিন}}, \emph{skrin}.) To download? (\textbf{\bn{ডাউনলোড} \bn{করা}}, \emph{ḍāunloḍ kôrā}.) A robot? (\textbf{\bn{রোবট}}, \emph{robôṭ}.) Automatic? (\textbf{\bn{স্বয়ংক্রিয়}}, \emph{svôyôngkriyô}.) A generator? (\textbf{\bn{জেনারেটর}}, \emph{jenāreṭôr}.) To rust? (\textbf{\bn{জং} \bn{ধরা}}, \emph{jông dhôrā}.) A motor? (\textbf{\bn{মোটর}}, \emph{moṭôr}.) Darkness? (\textbf{\bn{অন্ধকার}}, \emph{ôndhôkār}.) Bright? (\textbf{\bn{উজ্জ্বল}}, \emph{ujjvôl}.) A chandelier? (\textbf{\bn{ঝাড়বাতি}}, \emph{jhāṛbāti}.) To flicker? (\textbf{\bn{টিমটিম} \bn{করা}}, \emph{ṭimṭim kôrā}.) To go out (of a light)? (\textbf{\bn{নেভা}}, \emph{nebhā}.) A water tank? (\textbf{\bn{ট্যাংক}}, \emph{ṭængk}.) A hole right through? (\textbf{\bn{ফুটো}}, \emph{phuṭo}.) To overflow? (\textbf{\bn{উপচে} \bn{পড়া}}, \emph{upôche pôṛā}.) To bail out water? (\textbf{\bn{সেঁচা}}, \emph{sẽchā}.) To seep? (\textbf{\bn{চুঁইয়ে} \bn{পড়া}}, \emph{chũiye pôṛā}.) A small opening? (\textbf{\bn{ছিদ্র}}, \emph{chhidrô}.) Wall plaster? (\textbf{\bn{প্লাস্টার}}, \emph{plāsṭār}.) Glue? (\textbf{\bn{আঠা}}, \emph{āṭhā}.) To plaster? (\textbf{\bn{লেপা}}, \emph{lepā}.) To hammer in? (\textbf{\bn{ঠোকা}}, \emph{ṭhokā}.)
\end{tcolorbox}
